\section*{第 \ref{ch:b3:environmental-toxicology} 章——环境化学与毒理学}

\begin{solution}{exo:b3:environmental-toxicology:1}
同核双原子分子仅有一种振动，它使偶极矩保持为零，
氩没有振动，二者都不能吸收红外辐射。$\ce{CO2}$ 通过
弯曲振动（及反对称伸缩振动）会产生振荡的
偶极矩；弯曲振动恰好位于地球辐射的中段。
\end{solution}

\begin{solution}{exo:b3:environmental-toxicology:2}
$10^{0.1} = 1.26$：$[\ce{H+}]$ 上升 26\,\%。
\end{solution}

\begin{solution}{exo:b3:environmental-toxicology:3}
$(2.0 + 0.50)\ \unit{mmol/L} \times \qty{100.1}{mg/mmol} = \qty{250}{mg/L}$（以 $\ce{CaCO3}$ 计）。
\end{solution}

\begin{solution}{exo:b3:environmental-toxicology:4}
N：$0.60/14.0 = \qty{0.043}{mmol/L}$；P：$0.010/31.0 = \qty{3.2e-4}{mmol/L}$；N\,:\,P $= 133$，远高于 16，磷
首先耗尽，并限制生长。
\end{solution}

\begin{solution}{exo:b3:environmental-toxicology:5}
$L_0 = \mathrm{BOD_5}/(1 - \eu^{-0.23 \times 5}) = 170/0.683 = \qty{250}{mg/L}$.
\end{solution}

\begin{solution}{exo:b3:environmental-toxicology:6}
$\qty{4.40e-3}{L} \times \qty{0.0200}{mol/L} = \qty{8.80e-5}{mol}$，溶于 \qty{0.1000}{L} 中：碱度 $\qty{8.80e-4}{mol/L}$；以 $\ce{HCO3-}$ 计（\qty{61.0}{g/mol}），
\qty{53.7}{mg/L}。
\end{solution}

\begin{solution}{exo:b3:environmental-toxicology:7}
$K_d = 0.020 \times 300 = \qty{6.0}{L/kg}$。被吸附比例 $= K_dm_s/(K_dm_s + V_w) = 9.0/(9.0 + 0.30) = 0.97$：97\,\% 被吸附。
\end{solution}

\begin{solution}{exo:b3:environmental-toxicology:8}
苯：$\log\mathrm{BCF} = 0.85 \times 2.13 - 0.70 = 1.11$，BCF 约为 13。多氯联苯：$0.85 \times 6.67 - 0.70 = 4.97$，BCF 约为
\num{9e4}：即其在鱼体内的富集程度约为苯的七千倍；在这样的
$K_{ow}$ 下，该回归已接近其适用范围的极限。
\end{solution}

\begin{solution}{exo:b3:environmental-toxicology:9}
$\qty{192}{mg/kg} \times \qty{70}{kg} = \qty{13}{g}$，约相当于 150 杯。不同物种在吸收和代谢上存在差异，
且 $\mathrm{LD_{50}}$ 并非安全阈值；有害效应在远低于该值时就已开始出现。
\end{solution}

\begin{solution}{exo:b3:environmental-toxicology:10}
PNEC $= \qty{0.50}{mg/L}/1000 = \qty{0.50}{\micro g/L}$；$\mathrm{RQ} = 0.20/0.50 = 0.40$，小于 1：在此暴露水平下无需担忧。
\end{solution}

\begin{solution}{exo:b3:environmental-toxicology:11}
浮游生物/水：$\qty{0.015}{mg/kg}/\qty{2.0e-6}{mg/L} = \qty{7.5e3}{L/kg}$（生物富集）。鱼/浮游生物：$0.30/0.015 = 20$；
鸟/鱼：$6.0/0.30 = 20$（每一级均发生生物放大）；鸟/水：$\num{3e6}$。
\end{solution}

\begin{solution}{exo:b3:environmental-toxicology:12}
$\qty{1.0e9}{g}/\qty{64.1}{g/mol} = \qty{1.56e7}{mol}$ 的 $\ce{SO2}$，可生成相同摩尔数的 $\ce{H2SO4}$，质量为 $\qty{1.53e9}{g}$，约 \qty{1500}{t}。它
释放出 $\qty{3.12e7}{mol}$ 的 $\ce{H+}$；在 pH 4.0 时，$[\ce{H+}] = \qty{1.0e-4}{mol/L}$，即 $\qty{3.1e11}{L}$ 的雨水。
\end{solution}

\begin{solution}{pb:b3:environmental-toxicology:1}
\textbf{1.} 其在辛醇中的溶解度约为在水中的一万倍：具有疏水性，
因而会吸附于有机质上，并在脂肪中累积。

\textbf{2.} $K_{oc} = 0.41 \times 10^4 = \qty{4100}{L/kg}$; $K_{sw} = f_{oc}K_{oc} = 0.040 \times 4100 = \qty{164}{L/kg}$.

\textbf{3.} $K_{fw} = 0.048 \times 10^4 = \qty{480}{L/kg}$.

\textbf{4.} 否：平衡时空气中的浓度仅为水中浓度的 $10^{-4}$。

\textbf{5.} 这一中性的疏水分子溶解于有机质中，与
在辛醇中的情形相同；矿物表面是极性的，并被水覆盖。

\textbf{6.} 沉积物，其容量 $K_{sw}m_s$ 最大。

\textbf{7.} $V_w = \qty{1.0e9}{L}$; $K_{aw}V_a = \qty{1.0e8}{L}$; $K_{sw}m_s = \qty{1.64e10}{L}$; $K_{fw}m_f = \qty{4.8e6}{L}$.

\textbf{8.} $C_w = \qty{1.0e8}{mg}/\qty{1.75e10}{L} = \qty{5.7e-3}{mg/L}$, \qty{5.7}{\micro g/L}.

\textbf{9.} 水 \qty{5.7}{kg}，空气 \qty{0.57}{kg}，沉积物 \qty{94}{kg}，鱼 \qty{0.027}{kg}。

\textbf{10.} 空气 $\qty{5.7e-7}{mg/L}$；沉积物 \qty{0.94}{mg/kg}；鱼 \qty{2.7}{mg/kg}。

\textbf{11.} 各分区之间同时达到平衡，无降解、无
流入或流出，且各分区充分混合。

\textbf{12.} 在考虑降解并处于稳态（II 级）时，总量由
输入与损失速率决定；在考虑分区之间传输速率（III 级）时，
浓度不再符合平衡比例，而接收输入的
分区会被富集。

\textbf{13.} $k = \ln 2/\qty{60}{d} = \qty{0.0116}{d^{-1}}$.

\textbf{14.} $\eu^{-0.0116 \times 365} = 0.0145$：剩余 1.5\,\%。

\textbf{15.} $\qty{5.7}{\micro g/L} \times 0.0145 = \qty{0.083}{\micro g/L}$.

\textbf{16.} 在寒冷、缺氧的沉积物中，降解往往较慢，且被吸附的
农药缓慢释放回水体，使其在水体
已净化之后仍然存在。

\textbf{17.} 能抵抗水解、氧化和微生物作用的键（以芳香族
C--Cl 键为例）、水溶性低，以及吸附作用使其免受
降解者的作用。

\textbf{18.} PNEC $= \qty{50}{\micro g/L}/10 = \qty{5.0}{\micro g/L}$.

\textbf{19.} $\mathrm{RQ} = 5.7/5.0 = 1.1$.

\textbf{20.} 略高于 1：提示存在风险，评估须加以细化（采用更好的
毒性数据、实测浓度），或降低暴露水平。

\textbf{21.} \qty{2.7}{mg/kg}，几乎为 \qty{1.0}{mg/kg} 阈值的三倍：鱼
在浓度下降之前不宜食用。

\textbf{22.} 水、沉积物和鱼中的浓度随时间变化，并在若干采样点测定，
以检验模型并跟踪浓度的下降情况。

\textbf{23.} 减少农药施用量，避开降雨时段并远离岸边施用，
设置截留径流的缓冲带，或改用持久性较低的替代品。

\textbf{24.} 它将实测的无效应水平除以一个因子，以弥合少数
实验室物种与整个生态系统之间的差距；若因子取 10 而非 1000，则
PNEC 与结论都将相差一百倍。

\textbf{25.} 风险商 PEC/PNEC 约为 1.1，略高于 1。
\end{solution}
